\documentclass[aps,preprint]{revtex4}%
\usepackage{amssymb}
\usepackage{amsfonts}%
\usepackage{amsmath}%
\setcounter{MaxMatrixCols}{30}%
\usepackage{graphicx}
%TCIDATA{OutputFilter=latex2.dll}
%TCIDATA{Version=5.50.0.2960}
%TCIDATA{CSTFile=revtex4.cst}
%TCIDATA{Created=Wednesday, September 19, 2018 13:00:03}
%TCIDATA{LastRevised=Saturday, September 22, 2018 13:32:22}
%TCIDATA{<META NAME="GraphicsSave" CONTENT="32">}
%TCIDATA{<META NAME="SaveForMode" CONTENT="1">}
%TCIDATA{BibliographyScheme=Manual}
%TCIDATA{<META NAME="DocumentShell" CONTENT="Articles\SW\REVTeX 4">}
%BeginMSIPreambleData
\providecommand{\U}[1]{\protect\rule{.1in}{.1in}}
%EndMSIPreambleData
\newtheorem{theorem}{Theorem}
\newtheorem{acknowledgement}[theorem]{Acknowledgement}
\newtheorem{algorithm}[theorem]{Algorithm}
\newtheorem{axiom}[theorem]{Axiom}
\newtheorem{claim}[theorem]{Claim}
\newtheorem{conclusion}[theorem]{Conclusion}
\newtheorem{condition}[theorem]{Condition}
\newtheorem{conjecture}[theorem]{Conjecture}
\newtheorem{corollary}[theorem]{Corollary}
\newtheorem{criterion}[theorem]{Criterion}
\newtheorem{definition}[theorem]{Definition}
\newtheorem{example}[theorem]{Example}
\newtheorem{exercise}[theorem]{Exercise}
\newtheorem{lemma}[theorem]{Lemma}
\newtheorem{notation}[theorem]{Notation}
\newtheorem{problem}[theorem]{Problem}
\newtheorem{proposition}[theorem]{Proposition}
\newtheorem{remark}[theorem]{Remark}
\newtheorem{solution}[theorem]{Solution}
\newtheorem{summary}[theorem]{Summary}
\newenvironment{proof}[1][Proof]{\noindent\textbf{#1.} }{\ \rule{0.5em}{0.5em}}
\begin{document}
\preprint{HEP/123-qed}
\title[Short title for running header]{REV\TeX4}
\author{First Author}
\affiliation{My Institution}
\author{Second Author}
\affiliation{My Institution}
\author{Third Author}
\affiliation{Other Institution}
\keywords{one two three}
\pacs{PACS number}

\begin{abstract}
Shell document for REV\TeX{} 4.

\end{abstract}
\volumeyear{year}
\volumenumber{number}
\issuenumber{number}
\eid{identifier}
\date[Date text]{date}
\received[Received text]{date}

\revised[Revised text]{date}

\accepted[Accepted text]{date}

\published[Published text]{date}

\startpage{101}
\endpage{102}
\maketitle
\tableofcontents


\section{Thermal Time}

First, get the algebra $\mathcal{A}$ acting on a Hilbert space $\mathcal{H}$,
with a cyclic vector $\Phi_{\U{3c9} }\left(  I_{\mathcal{A}}\right)  $ (i.e.
such that $\Pi_{\omega}(\mathcal{A})\Phi_{\U{3c9} }\left(  I_{\mathcal{A}%
}\right)  $ is dense in $\mathcal{H}$ -- one way to get this is by the GNS
representation, so that the state $\U{3c9} $ just acts on an operator $A$ by
the expectation value at $\Phi_{\U{3c9} }\left(  I_{\mathcal{A}}\right)  $, as
above, so that the vector $\Phi_{\U{3c9} }\left(  I_{\mathcal{A}}\right)  $ is
standing in, in the Hilbert space picture, for the state $\U{3c9} $). Then one
can define an operator $S$ by the fact that, for any $A$ in $\mathcal{A}$, one
has%
\begin{equation}
S\left(  \Pi_{\U{3c9} }\left(  A\right)  \right)  \Phi_{\U{3c9} }\left(
I_{\mathcal{A}}\right)  =\Pi_{\U{3c9} }\left(  A\right)  ^{\dagger}%
\Phi_{\U{3c9} }\left(  I_{\mathcal{A}}\right)
\end{equation}


That is, $S$ acts like the conjugation operation on operators at
$\Phi_{\U{3c9} }\left(  I_{\mathcal{A}}\right)  $, which is enough to define
$S$ since $\Phi_{\U{3c9} }\left(  I_{\mathcal{A}}\right)  $ is cyclic. This
$S$ has a polar decomposition (analogous for operators to the polar form for
complex numbers) of
\begin{equation}
S=\mathcal{J\Delta}^{\frac{1}{2}},
\end{equation}
where$\mathcal{J}$ is antiunitary (this is conjugation, after all) and
$\mathcal{\Delta}$ is self-adjoint. We need the self-adjoint part, because the
Tomita flow is a one-parameter family of automorphisms given by:%
\begin{equation}
\widehat{\alpha\left(  t\right)  }\Pi_{\U{3c9} }\left(  A\right)
=\mathcal{\Delta}^{-it}\Pi_{\U{3c9} }\left(  A\right)  \mathcal{\Delta}^{it}.
\end{equation}


This leads to the definition of the funddamental Hamiltonian $\mathfrak{h}$
\begin{align}
\mathcal{\Delta}^{-it}  & =\exp\left\{  -i\mathfrak{h}t\right\}  \\
i\ln\left\{  \mathcal{\Delta}^{-i}\right\}    & =\mathfrak{h}\\
\mathfrak{h}  & \mathfrak{=}\ln\left\{  \mathcal{\Delta}\right\}  .
\end{align}


An important fact for Connes' classification of von Neumann algebras is that
the Tomita flow is basically unique -- that is, it's unique up to an inner
automorphism (i.e. a conjugation by some unitary operator -- so in particular,
if we're talking about a relativistic physical theory, a change of coordinates
giving a different $t$ parameter would be an example). So while there are
different flows, they're all \textquotedblleft essentially\textquotedblright%
\ the same. \emph{There's a unique notion of time flow if we reduce the
algebra }$\mathcal{A}$\emph{ to its cosets modulo inner automorphism i.e.}%
\begin{equation}
\operatorname*{Out}\left\{  \mathcal{A}\right\}  =\left.  \operatorname{Aut}%
\left\{  \mathcal{A}\right\}  \right/  \operatorname*{Inn}\left\{
\mathcal{A}\right\}  .
\end{equation}
Now, in some cases, the Tomita flow consists entirely of inner automorphisms,
and this reduction makes it disappear entirely (this happens in the
finite-dimensional case, for instance). But in the general case this doesn't
happen, and the Connes-Rovelli paper summarizes this by saying that von
Neumann algebras are \textquotedblleft intrinsically dynamic
objects\textquotedblright. So this is one interesting thing about the quantum
view of states: there is a somewhat canonical notion of dynamics present just
by virtue of the way states are described. In the classical world, this isn't
the case.

Now, Rovelli's \textquotedblleft Thermal Time\textquotedblright\ hypothesis
is, basically, that the notion of time is a state-dependent one: instead of an
independent variable, with respect to which other variables change, quantum
mechanics (per Rovelli) makes predictions about correlations between different
observed variables. More precisely, the hypothesis is that, given that we
observe the world in some state, the right notion of time should just be the
Tomita flow for that state. They claim that checking this for certain
cosmological models, like the Friedman model, they get the usual notion of
time flow. I have to admit, I have trouble grokking this idea as fundamental
physics, because it seems like it's implying that the universe (or any system
in it we look at) is always, a priori, in thermal equilibrium, which seems
wrong to me since it evidently isn't. The Friedman model does assume an
expanding universe in thermal equilibrium, but clearly we're not in exactly
that world. On the other hand, the Tomita flow is definitely there in the von
Neumann algebra view of quantum mechanics and states, so possibly I'm
misinterpreting the nature of the claim. Also, as applied to quantum gravity,
a \textquotedblleft state\textquotedblright\ perhaps should be read as a state
for the whole spacetime geometry of the universe -- which is presumably static
-- and then the apparent \textquotedblleft time change\textquotedblright%
\ would then be a result of the Tomita flow on operators describing actual
physical observables. But on this view, I'm not sure how to understand
\textquotedblleft thermal equilibrium\textquotedblright. So in the end, I
don't really know how to take the \textquotedblleft Thermal Time
Hypothesis\textquotedblright\ as physics.

In any case, the idea that the right notion of time should be state-dependent
does make some intuitive sense. The only physically, empirically accessible
referent for time is \textquotedblleft what a clock measures\textquotedblright%
: in other words, there is some chosen system which we refer to whenever we
say we're \textquotedblleft measuring time\textquotedblright. Different
choices of system (that is, different clocks) will give different readings
even if they happen to be moving together in an inertial frame -- atomic
clocks sitting side by side will still gradually drift out of sync. Even if
\textquotedblleft the system\textquotedblright\ means the whole universe, or
just the gravitational field, clearly the notion of time even in General
Relativity depends on the state of this system. If there is a
non-state-dependent \textquotedblleft god's-eye view\textquotedblright\ of
which variable is time, we don't have empirical access to it. So while I can't
really assess this idea confidently, it does seem to be getting at something important.


\end{document}